% ============================================================================
% paper9_v7_supplementary.tex --- Supplementary Material for paper9 v7
% ============================================================================
% Stand-alone document, included as Supplementary Information at submission.
% References are tables/sections cited from the main manuscript (paper9_v7_draft).
%
% Built 2026-05-23 from:
%   - tab:mpc-ablation lifted from paper9_v6.tex L267-290
%   - SI Tables S2-S4 generated by scripts/si_tables_v7.py from 175-cell sweep
%   - §S1 in-house DRL configuration documented from project memory
% ============================================================================

\documentclass[11pt]{article}
\usepackage[T1]{fontenc}
\usepackage[utf8]{inputenc}
\usepackage{amsmath,amssymb}
\usepackage{booktabs}
\usepackage{graphicx}
\usepackage{multirow}
\usepackage{caption}
\usepackage{hyperref}
\usepackage{natbib}
\usepackage[margin=1in]{geometry}
\bibliographystyle{plainnat}

% renumber tables/sections with S prefix
\renewcommand{\thetable}{S\arabic{table}}
\renewcommand{\thesection}{S\arabic{section}}
\renewcommand{\thefigure}{S\arabic{figure}}

\title{Supplementary Material for ``Model-based AI planning enables county-scale farmland consolidation in fragmented mountain landscapes''}

\author{Ning Zhou and Xiang Jing}

\date{}

\begin{document}
\maketitle
\sloppy  % allow looser line breaks for tech terms

% ============================================================================
\section{In-house model-free DRL baseline configuration}\label{si:s1-drl}
% ============================================================================

The Centralised PPO and Multi-Agent Reinforcement Learning (MARL) baselines reported in main-text Section~5 (and used as the fairness anchor for the contrastive-MPC results) were trained in-house under the following protocol.

\paragraph{Environment.} Both baselines act on the same Bishan County county-level environment (\texttt{CountyLevelEnv}) used by the contrastive MPC pipeline: 13 townships, $2{,}600$ spatial blocks aggregating $52{,}515$ cadastral parcels. The observation is the concatenation of $2{,}600$ blocks $\times$ 17 block-level features plus 12 global features. The action space is a categorical choice over the $2{,}600$ blocks, with the same MaskablePPO-style action mask described in main-text Methods.

\paragraph{Reward.} Identical to the main-text reward (Equation~1): $w_S = 4{,}000$, $w_C = 500$, $w_B = 1{,}500$, $b_B = 5$. No method-specific re-tuning was performed, so cross-method comparisons reflect the learning procedure rather than reward engineering.

\paragraph{Centralised PPO.} A single MaskablePPO agent with our \texttt{ParcelScoringPolicy} architecture (block-wise scoring MLP $[64, 32] \to 1$, global-feature value network). Training: 5 random seeds, $1.5 \times 10^6$ environment steps each, $\gamma = 0.995$, $\lambda = 0.95$, learning rate $1\times 10^{-3}$ linear-decayed, entropy coefficient $0.01$, clip range $0.2$, batch size $4{,}096$, mini-batch size $256$, 10 epochs per update. Wall-time: 8--12 GPU-hours per seed on a single A100. Cross-seed slope improvement: $-0.79\% \pm 0.36\%$, with $2/5$ seeds collapsing into degenerate trajectories.

\paragraph{Multi-Agent Reinforcement Learning (MARL).} A 13-agent variant in which each township is a separate MaskablePPO agent with shared parameters and a centralised value function over the global state, trained jointly. Same hyperparameters as Centralised PPO. 5 seeds, all stable. Cross-seed slope improvement: $-0.81\% \pm 0.10\%$ (CV $\approx 12\%$).

\paragraph{Why these are the relevant baselines.} Both baselines are task-faithful: same environment, same action space, same reward function, same evaluation protocol (5-episode deterministic rollout) as the contrastive MPC pipeline. The wall-time and seed-failure characteristics motivate the search for a less compute-intensive learned-model alternative; the slope-improvement magnitudes ($-0.8\%$ range) provide a non-trivial reference against which the contrastive-MPC result ($-1.289\% \pm 0.079\%$ on Bishan) is then compared.

% ============================================================================
\section{MPC optimisation ablation on the MSE-only ensemble (Supplementary Table S1)}\label{si:s2-mpc-ablation}
% ============================================================================

Table~\ref{tab:si-mpc-ablation} below reports the full MPC optimisation ablation cited in main-text Section~3 (paragraph beginning ``Four independent operational failures''). All variants use the same MSE-only world-model ensemble (Bishan District, 5-episode evaluation) and modify a single dimension of the planner. None exceeds the random-sampling baseline; the catastrophic failure of greedy continuation under the MSE-only ensemble (slope $+0.029\%$, wrong direction) is the variant that contrastive training subsequently rescues into the strongest variant tested (main-text Table~3).

\begin{table}[htbp]
\centering
\small
\caption{MPC optimisation ablation across four dimensions, on the MSE-only world-model ensemble (Bishan District). All variants start from the optimal $H{=}5$, $K{=}50$ random-reward configuration and modify one dimension. None exceeds the random-sampling baseline; aggressive exploitation (greedy continuation) actively reverses the slope direction. The diagnostic shows that the failure is structural to the learned model rather than to the search procedure.}
\label{tab:si-mpc-ablation}
\resizebox{\textwidth}{!}{%
\begin{tabular}{llccc}
\toprule
Dimension & Configuration & Slope \% & Reward & Diagnosis \\
\midrule
\multirow{2}{*}{Continuation} & Random (baseline) & $-0.209 \pm 0.008$ & $-18.52$ & --- \\
 & Greedy & $+0.029 \pm 0.005$ & $-95.29$ & catastrophic \\
\midrule
\multirow{2}{*}{Scoring} & Reward (baseline) & $-0.209 \pm 0.008$ & $-18.52$ & --- \\
 & Slope-only & $-0.001$ & $-98.95$ & zero discrimination \\
\midrule
\multirow{2}{*}{Terminal value} & No value (baseline) & $-0.209 \pm 0.008$ & $-18.52$ & --- \\
 & Value-guided ($v{=}1.0$) & $-0.172 \pm 0.042$ & $-27.96$ & no improvement \\
\midrule
\multirow{3}{*}{Search budget} & $K{=}50$ (baseline) & $-0.209 \pm 0.008$ & $-18.52$ & --- \\
 & $K{=}100$ & $-0.167 \pm 0.053$ & $-8.94$ & no improvement \\
 & $K{=}50$, $n_r{=}3$ & $-0.171 \pm 0.046$ & $-28.33$ & no improvement \\
\bottomrule
\end{tabular}%
}
\end{table}

% ============================================================================
\section{Synthetic-benchmark secondary metrics (Tables S2--S4)}\label{si:s3-bench-secondary}
% ============================================================================

Main-text Table~5 (\texttt{tab:bench-slope}) reports the slope-improvement profile across the seven synthetic landscapes. The three secondary metrics referenced in main-text Section~6 ($\Delta$contiguity, $\Delta$\textit{baimu fang} count, and $\Delta$\textit{baimu fang} area, all under the same 50-pair swap budget and $n{=}5$ seeds per cell) are tabulated here. Together with main-text Table~5 these summarise the full $5 \times 7 \times 5 = 175$-cell sweep. Aggregation script: \texttt{scripts/si\_tables\_v7.py} in the released code repository.

\begin{table}[htbp]
\centering
\small
\caption{Synthetic-benchmark contiguity change ($\Delta$Contiguity, absolute units). Mean $\pm$ standard deviation across $n{=}5$ seeds. Higher (more positive) is better; per-row best in bold.}
\label{tab:si-bench-cont}
\resizebox{\textwidth}{!}{%
\begin{tabular}{lccccc}
\toprule
Preset ($n_{\text{blocks}}$) & Random & Greedy & GA & PPO & MPC \\
\midrule
\texttt{bishan\_clone} (2{,}600) & $0.0016{\pm}0.0007$ & $0.0009{\pm}0.0006$ & $\mathbf{0.0253{\pm}0.0023}$ & $0.0027{\pm}0.0005$ & $0.0032{\pm}0.0009$ \\
\texttt{neijiang\_clone} (3{,}700) & $0.0014{\pm}0.0004$ & $0.0012{\pm}0.0001$ & $\mathbf{0.0159{\pm}0.0005}$ & $0.0015{\pm}0.0004$ & $0.0023{\pm}0.0005$ \\
\texttt{plain\_small\_cons} (800) & $0.0057{\pm}0.0023$ & $0.0036{\pm}0.0010$ & $0.0018{\pm}0.0057$ & $\mathbf{0.0069{\pm}0.0014}$ & $0.0068{\pm}0.0015$ \\
\texttt{plain\_large\_cons} (4{,}500) & $0.0010{\pm}0.0001$ & $0.0008{\pm}0.0002$ & $0.0010{\pm}0.0003$ & $\mathbf{0.0014{\pm}0.0002}$ & $0.0012{\pm}0.0003$ \\
\texttt{plain\_medium\_frag} (2{,}600) & $0.0018{\pm}0.0004$ & $0.0018{\pm}0.0005$ & $0.0005{\pm}0.0006$ & $0.0022{\pm}0.0003$ & $\mathbf{0.0026{\pm}0.0004}$ \\
\texttt{mixed\_medium\_frag} (2{,}600) & $0.0022{\pm}0.0006$ & $0.0018{\pm}0.0003$ & $\mathbf{0.0059{\pm}0.0016}$ & $0.0031{\pm}0.0007$ & $0.0032{\pm}0.0008$ \\
\texttt{hilly\_small\_cons} (800) & $0.0058{\pm}0.0020$ & $0.0067{\pm}0.0017$ & $\mathbf{0.0734{\pm}0.0055}$ & $0.0092{\pm}0.0020$ & $0.0095{\pm}0.0013$ \\
\bottomrule
\end{tabular}%
}
\end{table}

\begin{table}[htbp]
\centering
\small
\caption{Synthetic-benchmark large-patch count change ($\Delta$\textit{baimu fang} count, $\geq 6.67$\,ha qualifying patches gained per episode). Mean $\pm$ standard deviation across $n{=}5$ seeds. Higher is better; per-row best in bold.}
\label{tab:si-bench-baimuct}
\resizebox{\textwidth}{!}{%
\begin{tabular}{lccccc}
\toprule
Preset ($n_{\text{blocks}}$) & Random & Greedy & GA & PPO & MPC \\
\midrule
\texttt{bishan\_clone} (2{,}600) & $0.0{\pm}2.2$ & $-0.2{\pm}2.6$ & $-3.8{\pm}4.4$ & $\mathbf{9.4{\pm}2.6}$ & $9.2{\pm}2.3$ \\
\texttt{neijiang\_clone} (3{,}700) & $0.6{\pm}0.9$ & $-1.4{\pm}1.5$ & $-4.4{\pm}3.2$ & $\mathbf{6.4{\pm}2.1}$ & $4.6{\pm}3.4$ \\
\texttt{plain\_small\_cons} (800) & $0.0{\pm}0.0$ & $0.2{\pm}1.1$ & $\mathbf{0.8{\pm}2.0}$ & $-0.4{\pm}2.5$ & $0.4{\pm}2.1$ \\
\texttt{plain\_large\_cons} (4{,}500) & $0.0{\pm}0.0$ & $0.0{\pm}0.0$ & $0.2{\pm}0.4$ & $\mathbf{0.8{\pm}0.8}$ & $0.0{\pm}0.7$ \\
\texttt{plain\_medium\_frag} (2{,}600) & $0.6{\pm}0.9$ & $-0.6{\pm}1.3$ & $1.4{\pm}1.8$ & $\mathbf{1.6{\pm}1.1}$ & $0.4{\pm}1.3$ \\
\texttt{mixed\_medium\_frag} (2{,}600) & $-0.2{\pm}1.3$ & $0.0{\pm}1.2$ & $\mathbf{1.2{\pm}2.2}$ & $-1.0{\pm}2.2$ & $-2.4{\pm}1.1$ \\
\texttt{hilly\_small\_cons} (800) & $-1.0{\pm}0.7$ & $-0.4{\pm}0.9$ & $-3.2{\pm}2.2$ & $3.4{\pm}1.7$ & $\mathbf{9.8{\pm}1.5}$ \\
\bottomrule
\end{tabular}%
}
\end{table}

\begin{table}[htbp]
\centering
\small
\caption{Synthetic-benchmark large-patch area change ($\Delta$\textit{baimu fang} area, ha). Mean $\pm$ standard deviation across $n{=}5$ seeds. The sign of this metric is morphology-dependent (see main-text Section~5 discussion of the Bishan vs.\ Neijiang reversal); per-row best (most positive) in bold for reference.}
\label{tab:si-bench-baimuha}
\resizebox{\textwidth}{!}{%
\begin{tabular}{lccccc}
\toprule
Preset ($n_{\text{blocks}}$) & Random & Greedy & GA & PPO & MPC \\
\midrule
\texttt{bishan\_clone} (2{,}600) & $9.8{\pm}11.1$ & $11.9{\pm}15.0$ & $155.3{\pm}40.3$ & $134.7{\pm}26.1$ & $\mathbf{177.5{\pm}31.7}$ \\
\texttt{neijiang\_clone} (3{,}700) & $15.2{\pm}9.2$ & $27.8{\pm}16.4$ & $\mathbf{190.8{\pm}33.2}$ & $122.4{\pm}44.8$ & $167.6{\pm}23.0$ \\
\texttt{plain\_small\_cons} (800) & $20.2{\pm}8.9$ & $14.7{\pm}16.2$ & $11.5{\pm}19.5$ & $19.0{\pm}25.3$ & $\mathbf{30.5{\pm}17.5}$ \\
\texttt{plain\_large\_cons} (4{,}500) & $9.2{\pm}8.9$ & $39.7{\pm}13.3$ & $12.8{\pm}13.4$ & $40.1{\pm}55.3$ & $\mathbf{40.3{\pm}15.3}$ \\
\texttt{plain\_medium\_frag} (2{,}600) & $22.8{\pm}10.5$ & $38.5{\pm}22.8$ & $34.2{\pm}34.3$ & $\mathbf{54.0{\pm}40.1}$ & $52.8{\pm}26.8$ \\
\texttt{mixed\_medium\_frag} (2{,}600) & $14.2{\pm}7.3$ & $28.4{\pm}13.6$ & $60.7{\pm}24.8$ & $\mathbf{106.8{\pm}24.2}$ & $102.1{\pm}13.1$ \\
\texttt{hilly\_small\_cons} (800) & $21.6{\pm}14.2$ & $4.6{\pm}13.9$ & $\mathbf{167.7{\pm}41.2}$ & $100.3{\pm}30.3$ & $131.9{\pm}43.5$ \\
\bottomrule
\end{tabular}%
}
\end{table}


% ============================================================================
\section{Toolchain pre-flight checks}\label{si:s4-toolchain}
% ============================================================================

Tool 5 of the ArcGIS Pro toolchain (main-text Section~7) validates the user environment before any planning step is allowed to run. The complete list of checks is:

\begin{enumerate}
\item \textbf{Python interpreter}: ArcGIS Pro 3.6+ \texttt{arcgispro-py3} or compatible Python 3.10+ environment.
\item \textbf{Spatial Analyst extension}: licence checked out (required by Tool~1 for slope derivation).
\item \textbf{Required packages}: \texttt{gymnasium} ($\geq 0.29$), \texttt{libpysal} ($\geq 4.7$), \texttt{onnxruntime} ($\geq 1.16$), \texttt{numpy}, \texttt{pandas}, \texttt{scipy}.
\item \textbf{ONNX-ensemble shape compatibility}: if a pre-trained ensemble is provided, the baked-in $N_{\text{blocks}}$ matches the prepared dataset.
\item \textbf{Coordinate reference system}: the input shapefile's CRS resolves to a metric projection (default \texttt{EPSG:32648}; user-overridable).
\item \textbf{Disk space}: at least 5\,GB free in the prepared-data directory.
\end{enumerate}

Tool~5 prints an \texttt{[OK]} or \texttt{[FAIL]} verdict for each item and refuses to proceed with downstream tools if any required item fails. The 70-second smoke test on the bundled 2-township subset of Neijiang Dongxing District exercises every check end to end and is the recommended onboarding action for new users of the toolchain.

% ============================================================================
\section{Data and code release}\label{si:s5-release}
% ============================================================================

The synthetic benchmark, the pre-trained Bishan ensemble, the ArcGIS Pro toolbox source, and the aggregation/sweep scripts referenced throughout this paper are released under the licences indicated in the main-text Data and Code Availability statements. Repository structure at the public release:

\begin{itemize}
\item \texttt{LandUseOptimization\_P9.pyt}: ArcGIS Pro Python Toolbox manifest.
\item \texttt{farmland\_mpc/}: pure-Python core (loss functions, ensemble training, MPC planner, ONNX export).
\item \texttt{benchmark/}: synthetic generator, presets, baseline runners, sweep manifest, and the 175-cell results.
\item \texttt{docs/BENCHMARK\_RESULTS.md}: per-(preset, method) aggregated results derived from \texttt{benchmark/results/}.
\item \texttt{scripts/aggregate\_results.py}, \texttt{scripts/si\_tables\_v7.py}: aggregation and SI-table generators.
\item \texttt{environment.yml}, Colab demo notebook, integration smoke test.
\end{itemize}

% ============================================================================
\section{Verification stack: full per-episode data (Tables S5--S8)}\label{si:s5-verification}
% ============================================================================

The main-text Methods sections (\texttt{Independent GIS-grounded verification}, \texttt{Direct measurement of ensemble 1-step prediction accuracy}, and \texttt{Replacing the ensemble with the true environment}) summarise four of the five independent verification layers underwriting the Bishan headline result (the fifth, Neijiang cross-county replication, is in main-text Section~5). Tables~\ref{tab:si-validate}--\ref{tab:si-trueenv} give the per-replicate, per-channel numbers for each verification layer so external readers can audit them without rerunning the released scripts.

\paragraph{Layer (i): GIS-grounded recomputation.} Table~\ref{tab:si-validate} compares the planner-reported deltas (from \texttt{mpc\_summary.json}) against an independent recomputation that reads only the optimised DLTB shapefile, joins per-parcel slope from the prepared DEM raster, builds Queen contiguity via \texttt{libpysal}, and applies the same metric definitions as the main-text Methods. The recomputation script (\texttt{validate\_optimized\_shp.py}) shares no source with the optimisation pipeline.

\begin{table}[htbp]
\centering
\small
\caption{Independent GIS recomputation versus planner-reported deltas, Bishan run\_1, $53{,}086$ swappable parcels. The recomputation invokes no learned model. Differences are at floating-point reduction noise level.}
\label{tab:si-validate}
\begin{tabular}{lrrr}
\toprule
Quantity & Recomputed (independent) & Planner-reported & $|\Delta|$ \\
\midrule
$\Delta$ slope (\%) & $-2.0000$ & $-2.0006$ & $6 \times 10^{-4}$ pp \\
$\Delta$ contiguity & $+0.012739$ & $+0.012748$ & $8 \times 10^{-6}$ \\
$\Delta$ baimu count & $+8$ & $+8$ & $0$ (exact) \\
$\Delta$ baimu area (ha) & $-473.2950$ & $-473.2950$ & $< 10^{-4}$ ha \\
farm$\rightarrow$forest swaps & $428$ & $428$ & $0$ (exact) \\
forest$\rightarrow$farm swaps & $428$ & $428$ & $0$ (exact) \\
\bottomrule
\end{tabular}
\end{table}

\paragraph{Layer (ii): ensemble member-subset ablation.} Table~\ref{tab:si-subset} reports the three replicates that constitute the variance estimator used in the main-text Methods caveat on five-seed determinism. Each replicate runs the full MPC pipeline (H=5, K=50, greedy continuation, scoring=reward) with a different 2-of-3 ensemble member subset. The full ensemble's result lies inside the mean $\pm$~sd of the subset distribution on every reported metric.

\begin{table}[htbp]
\centering
\small
\caption{Per-replicate results for the member-subset ablation. The three replicates exhaust C(3,2)=3 distinct subsets. Full-ensemble values reproduced from main-text Table~4 / \texttt{mpc\_summary.json} for comparison; they fall inside the subset mean$\pm$sd on every metric.}
\label{tab:si-subset}
\begin{tabular}{lcccc}
\toprule
Replicate (members) & $\Delta$ slope (\%) & $\Delta$ cont & $\Delta$ baimu count & $\Delta$ baimu (ha) \\
\midrule
\{0, 1\} & $-2.0222$ & $+0.01361$ & $+5$ & $-475.02$ \\
\{0, 2\} & $-1.9706$ & $+0.01437$ & $+5$ & $-443.12$ \\
\{1, 2\} & $-2.0024$ & $+0.01128$ & $+6$ & $-505.39$ \\
\midrule
Mean $\pm$ sd (n=3) & $-1.9984 \pm 0.0260$ & $+0.01309 \pm 0.00161$ & $+5.33 \pm 0.58$ & $-474.51 \pm 31.14$ \\
Full ensemble \{0,1,2\} & $-2.0006$ & $+0.01275$ & $+8$ & $-473.30$ \\
\bottomrule
\end{tabular}
\end{table}

\paragraph{Layer (iii): direct ensemble 1-step accuracy.} Table~\ref{tab:si-mae} reports the per-channel prediction error of the deployed 3-member ensemble against the true environment, accumulated over the 100 in-distribution steps that reproduce episode~0 of \texttt{mpc\_summary.json}. The reward head is strongly under-fit but Spearman rank correlation between predicted and true reward is positive, and slope-channel error is sub-variance.

\begin{table}[htbp]
\centering
\small
\caption{Ensemble 1-step prediction error over the same 100-step trajectory MPC visits (Bishan, seed 0). True-reward statistics are computed from the actual env.step return. Spearman rank correlation between predicted and true reward is $+0.216$. Pearson correlation between ensemble disagreement $r_{\text{std}}$ and absolute prediction error is $+0.077$, i.e.\ disagreement does not flag where members are wrong.}
\label{tab:si-mae}
\begin{tabular}{lr}
\toprule
Quantity & Value \\
\midrule
Reward MAE & $0.818$ \\
Reward RMSE & $1.540$ \\
Reward median $|\text{err}|$ & $0.170$ \\
Reward p90 $|\text{err}|$ & $3.148$ \\
Reward p99 $|\text{err}|$ & $5.075$ \\
$r_{\text{true}}$ range / std & $[-3.04, +6.36]$ / $1.397$ \\
$r_{\text{pred}}$ range / std & $[+0.43, +0.67]$ / $0.069$ \\
Naive constant-mean baseline MAE & $0.828$ \\
Spearman $(r_{\text{pred}}, r_{\text{true}})$ & $+0.216$ \\
Ensemble $r_{\text{std}}$ mean / median / max & $0.027$ / $0.030$ / $0.075$ \\
Pearson $(r_{\text{std}}, |\text{err}|)$ & $+0.077$ \\
$gf[4]$ slope-improvement MAE & $0.00238$ (channel std $0.00521$) \\
$gf[6]$ baimu-count-norm MAE & $0.00270$ (channel std $0.00817$) \\
\bottomrule
\end{tabular}
\end{table}

\paragraph{Layer (iv): true-environment MPC counter-factual.} Table~\ref{tab:si-trueenv} compares the ensemble-MPC headline trajectory against an ablation that replaces every ensemble call with the actual env.step, under matched H/K/$\gamma$ and a single seed. Implementation concessions (random rather than greedy continuation; stage-1 subsample of 200 of $\sim$2{,}000 valid actions) are forced by the cost of incremental simulation. Despite using a perfect dynamics model, the resulting trajectory is strictly worse on all three objectives.

\begin{table}[htbp]
\centering
\small
\caption{True-environment MPC vs.\ deployed ensemble MPC, Bishan, 100 outer steps, seed 0. ``True-env'' replaces every ensemble call with actual env.step; otherwise H=5, K=50, $\gamma$=0.99 unchanged. Implementation concessions stated in caption: random (rather than greedy) continuation, stage-1 subsample of 200 valid actions per step.}
\label{tab:si-trueenv}
\begin{tabular}{lccc}
\toprule
Quantity & Ensemble MPC (run\_1) & True-env MPC & $\Delta$ (true $-$ ensemble) \\
\midrule
$\Delta$ slope (\%) & $-2.0006$ & $-1.8187$ & $+0.182$ pp (worse) \\
$\Delta$ contiguity & $+0.01275$ & $+0.00865$ & $-0.0041$ (worse) \\
$\Delta$ baimu count & $+8$ & $+2$ & $-6$ (worse) \\
$\Delta$ baimu area (ha) & $-473.30$ & $-565.10$ & $-91.8$ ha more loss \\
Wall time (excl.\ env build) & $1{,}448$ s & $358$ s & true-env 4$\times$ faster \\
Stage-1 coverage & all valid actions (batched) & 200-action random subset & — \\
Stage-2 continuation & greedy & random & — \\
\bottomrule
\end{tabular}
\end{table}

Per-step trajectory matching between layer~(iii) and the deployed pipeline is byte-identical at every monitored step (10, 20, \ldots, 100) of episode 0; this is the most stringent reproducibility check the verification stack supports and is therefore the canonical evidence that the deployed ensemble's planning is deterministic given a fixed prepared dataset and config.

\end{document}
